\chapter{Einführung}\label{kap:einfuehrung}
\index{Raytatouille Fusion}

Dieses Handbuch beschreibt \produkt{} in der Version \version. Es richtet sich an
Anwenderinnen und Anwender aus Optik und Physik, die Geometrie modellieren oder einlesen und
darauf optische Berechnungen ausführen wollen -- Strahlverfolgung, Wellenoptik, Beugung an
Gittern, Lichttransport in Szenen und Scheinwerfern. Programmierkenntnisse sind für die
Arbeit mit dem Programm nicht nötig; nur das Einrichten der Python-Umgebung für die Plugins
verlangt einige Schritte im Dateisystem, die in \autoref{kap:installation} genau beschrieben
sind.

\section{Was ist Raytatouille Fusion?}\label{sec:was-ist}

\produkt{} ist ein CAD-Programm für den Schreibtisch, das zwei Aufgaben verbindet:

\begin{enumerate}
  \item \textbf{Geometrie bearbeiten.} Das Programm ist die Bedienoberfläche
    (\enquote{Frontend}) für den NURBS-Kern \textbf{Nurbforge}\index{Nurbforge}. Alle Kurven,
    Flächen und Körper werden darin \emph{exakt} als NURBS (nicht-uniforme rationale
    B-Splines) gespeichert: mit Grad, Knotenvektor, Kontrollpunkten und Gewichten. Zur Anzeige
    werden sie in Dreiecke zerlegt (tesselliert); diese Näherung dient nur dem Bildschirm und
    dem STL-Export, gerechnet und gespeichert wird immer mit der exakten Geometrie.
  \item \textbf{Optisch rechnen.} Über Plugins werden Rechenkerne angebunden. Sie lesen das
    Modell, rechnen und zeigen ihre Ergebnisse direkt am Modell: Strahlen als Linien,
    Felder und Bestrahlungsstärken farbig auf den Flächen, Bilder in einem eigenen Fenster,
    Kurven in einem Diagramm, Tabellen in einer Tabellenansicht.
\end{enumerate}

Mit dem Programm ausgeliefert werden fünf Plugins, die auf die folgenden Rechenkerne
zugreifen:

\begin{referenz}{>{\raggedright\arraybackslash}p{3.2cm}X}
\toprule
\textbf{Rechenkern} & \textbf{Aufgabe} \\
\midrule
\endhead
\textbf{Raytatouille}\index{Raytatouille} & Strahlverfolgung durch optische Systeme
  (Linsen, Spiegel, Strahlteiler), Bestrahlungsstärke auf der Bildfläche, Optimierung. Wird
  über das Plugin \plugin{Raytracer} angesprochen (\autoref{kap:raytracer}). \\
\textbf{RTT-Fusion}\index{RTT-Fusion} & Die Verbindung von Strahlen- und Wellenoptik in einem
  Knoten-Workflow: Systeme werden in Blöcke zerlegt, zwischen denen Licht als Strahlen oder
  als Feld weitergegeben wird. Das Plugin \plugin{Workflow} ist die Rechenmaschine dieses
  Workflows (\autoref{kap:workflow}). \\
\textbf{Fieldtracing}\index{Fieldtracing} & Feldausbreitung (z.\,B.\ Debye,
  Rayleigh--Sommerfeld) für die Verbindungen des Workflows, die Felder statt Strahlen
  übertragen; wird vom Raytracer und vom Workflow mitbenutzt. \\
\textbf{FourierGrating}\index{FourierGrating} & Beugung an periodischen Strukturen (RCWA,
  differentielle Methode), Wirkungsgrade je Beugungsordnung. Plugin \plugin{FourierGrating}
  (\autoref{kap:fouriergrating}). \\
\textbf{Path Tracer}\index{Path Tracer} & Lichttransport auf dem CAD-Modell selbst:
  Objekte werden zu Lichtquellen, die übrigen erhalten Materialien; Ergebnis sind
  Bestrahlungsstärken auf den Flächen und Kamerabilder. Plugin \plugin{Path Tracer}
  (\autoref{kap:pathtracer}); dasselbe Rechenpaket nutzt auch das Plugin
  \plugin{Automotive Lighting} für Scheinwerfer und Signalleuchten
  (\autoref{kap:automotive}). \\
\bottomrule
\end{referenz}

Die Rechenkerne selbst sind Python-Pakete (\code{raytatouille}, \code{rtt-fusion},
\code{photonpower-fieldtracing}, \code{fouriergrating}, \code{pathtracer}). Sie werden nicht mit
dem Programm installiert, sondern in eine Python-Umgebung auf Ihrem Rechner; das Programm
richtet diese Umgebung auf Knopfdruck ein (\autoref{sec:python-umgebung}). Die Plugins laufen
jeweils in einem eigenen Python-Prozess. Stürzt eine Rechnung ab oder hängt sie, bleibt das
CAD-Programm bedienbar; eine laufende Rechnung lässt sich jederzeit abbrechen.

\begin{hinweis}
\produkt{} hieß bis November 2026 \enquote{Nurbforge CAD}. Ältere Dokumente mit der Endung
\datei{.nfcad} lassen sich weiterhin öffnen, und die Einstellungen von Nurbforge CAD werden
beim ersten Start übernommen (\autoref{sec:einstellungen}). An einigen Stellen ist der alte
Name noch sichtbar, etwa im Namen des Python-Pakets \code{nfcad} für Plugins oder in den
Umgebungsvariablen \code{NFCAD\_\dots}.
\end{hinweis}

\subsection*{Was das Programm kann}

\begin{itemize}
  \item STEP-Dateien und eigene Dokumente (\datei{.rttf}) öffnen und in vier Ansichten
    darstellen -- als Drahtmodell, mit flachen oder mit beleuchteten Flächen
    (\autoref{kap:hauptfenster}, \autoref{kap:navigation}).
  \item Objekte, einzelne Flächen, Kanten und Eckpunkte auswählen; ihre Eigenschaften (Grad,
    Knoten, Länge, Fläche, Volumen) ablesen (\autoref{kap:navigation},
    \autoref{kap:organisieren}).
  \item Mit einem \enquote{magnetischen} Cursor zeichnen, der auf Endpunkte, Mittelpunkte,
    Zentren, Schnittpunkte, Flächenecken und -kanten einrastet; Koordinaten eintippen
    (\autoref{kap:modellieren}).
  \item Kurven, Flächen (Extrusion, Rotation, Loft, Sweep, Coons, Gordon, Offsets,
    Verrundungen) und Körper (Quader, Kugel, Zylinder, Kegel, Torus, boolesche Operationen)
    erzeugen, transformieren und bearbeiten (\autoref{kap:modellieren}).
  \item Modelle mit Ebenen, Gruppen und freien Attributen ordnen (\autoref{kap:organisieren}).
  \item STEP (AP214) und STL schreiben (\autoref{kap:dateien}).
  \item Optische Rechnungen über Plugins ausführen und die Ergebnisse mit dem Dokument
    speichern (\autoref{kap:plugins} bis \autoref{kap:automotive}).
\end{itemize}

Das Programm läuft unter macOS (ab Version~14, Apple Silicon und Intel), Windows~11 und
AlmaLinux bzw.\ Red Hat Enterprise Linux~10. Es zeichnet mit der jeweils nativen
Grafikschnittstelle (Metal, Direct3D, OpenGL oder Vulkan).

\section{Grundbegriffe}\label{sec:begriffe}

Die folgenden Begriffe werden im ganzen Handbuch in dieser Bedeutung verwendet. Weitere
Fachbegriffe erklärt das Glossar (\autoref{kap:glossar}).

\begin{description}
  \item[Dokument]\index{Dokument} Alles, was Sie in einem Fenster bearbeiten: die Geometrie
    mit Ebenen, Gruppen und Attributen, die Parameter der Plugins, ein etwaiger
    Knoten-Workflow und die gespeicherten Ergebnisse. Ein Dokument wird als eine Datei mit
    der Endung \datei{.rttf} gespeichert (\autoref{kap:dateien}). Das Programm hat immer genau
    ein Dokument geöffnet; ein neues oder geöffnetes Dokument ersetzt das bisherige.
  \item[Objekt (Entität)]\index{Objekt}\index{Entität} Ein einzelnes Geometrieelement des
    Dokuments. Es gibt fünf Arten, die im Modellbaum in der Spalte \gui{Type} stehen:
    \begin{itemize}
      \item \gui{Point} -- ein Punkt,
      \item \gui{Curve} -- eine NURBS-Kurve (Linie, Polylinie, Kreis, Bogen, Spline \dots),
      \item \gui{Surface} -- eine untrimmte NURBS-Fläche,
      \item \gui{Trimmed surface} -- eine NURBS-Fläche, die durch Randkurven im
        Parameterraum beschnitten ist (z.\,B.\ mit Loch),
      \item \gui{Shell} -- eine Hülle aus mehreren zusammengenähten getrimmten Flächen
        (B-Rep). Geschlossene Hüllen sind Volumenkörper; STEP-Dateien werden als Hüllen
        eingelesen.
    \end{itemize}
    Die Begriffe Objekt und Entität sind gleichbedeutend; die Oberfläche spricht von
    \enquote{object} oder \enquote{entity}. Jedes Objekt hat eine eindeutige Nummer (\gui{Id}),
    optional einen Namen, gehört zu genau einer Ebene und höchstens einer Gruppe.
  \item[Fläche (Face)]\index{Fläche} Eine einzelne getrimmte Fläche innerhalb einer Hülle.
    Flächen sind \emph{Teilobjekte}: Sie lassen sich auswählen und ihre Eigenschaften ablesen,
    aber nicht für sich löschen.
  \item[Kante (Edge)]\index{Kante} Eine Randkurve einer Hülle, an der eine oder zwei Flächen
    zusammenstoßen; ebenfalls ein Teilobjekt.
  \item[Eckpunkt (Vertex)]\index{Eckpunkt} Ein Punkt, an dem Kanten einer Hülle enden;
    Teilobjekt.
  \item[Ebene (Layer)]\index{Ebene}\index{Layer} Eine Zusammenfassung von Objekten mit
    gemeinsamem Namen, gemeinsamer Farbe und gemeinsamem Sichtbarkeitsschalter; eine Ebene
    lässt sich sperren. Jedes Dokument hat mindestens die Ebene \gui{Default}; neue Objekte
    entstehen auf der \emph{aktuellen} Ebene. Nicht zu verwechseln mit der
    \emph{Konstruktionsebene} einer Ansicht (\autoref{sec:konstruktionsebene}), in der
    gezeichnet wird.
  \item[Gruppe]\index{Gruppe} Ein Knoten im Modellbaum, der Objekte und weitere Gruppen
    enthält. Gruppen dienen der Ordnung (ein Objektiv, eine Baugruppe aus einer STEP-Datei)
    und lassen sich gemeinsam ein- und ausblenden (\autoref{sec:gruppen}).
  \item[Attribut]\index{Attribut} Ein frei wählbares Schlüssel-Wert-Paar an einem Objekt,
    z.\,B.\ \code{sim.inactive = true}. Attribute werden mit dem Dokument gespeichert; Plugins
    lesen und schreiben sie (\autoref{sec:attribute}).
  \item[Plugin]\index{Plugin} Eine Erweiterung in Python, die einen Rechenkern anbindet. Ein
    Plugin bringt eine Gruppe von Knöpfen in der Werkzeugspalte, eine Seite im Dock
    \gui{Plugin} mit seinen Parametern und Aktionen und ein Untermenü im Menü \menu{Plugins}
    mit (\autoref{kap:plugins}).
  \item[Aktion]\index{Aktion} Eine Rechnung oder ein Arbeitsschritt eines Plugins, etwa
    \enquote{Trace rays} oder \enquote{Camera image}. Eine Aktion läuft als \emph{Job} im
    Hintergrund; der Fortschritt steht in der Statusleiste. Alles, was eine Aktion am Modell
    ändert, ist \emph{ein} Rückgängig-Schritt.
  \item[Workflow]\index{Workflow} Der Knoten-Workflow: ein Graph aus Blöcken (Quellen,
    optische Elemente, Detektoren, Anzeigen) und Verbindungen, der ein optisches System
    beschreibt und blockweise berechnet wird. Er gehört zum Dokument und wird im Fenster
    \gui{Node workflow} bearbeitet (\autoref{kap:workflow}).
  \item[Ergebnis]\index{Ergebnis} Was eine Plugin-Rechnung liefert: Strahlen (Linien), Felder
    auf Flächen, Bilder, Kurven, Tabellen oder Balkendiagramme. Ergebnisse einer Rechnung
    bilden einen \emph{Lauf} (\enquote{run}) im Dock \gui{Results}. Ergebnisse sind
    Anzeigedaten, keine Geometrie -- Strahlen lassen sich aber in CAD-Kurven umwandeln
    (\autoref{sec:dock-results}). Die neuesten Ergebnisse der Anzeigeblöcke des Workflows
    werden mit dem Dokument gespeichert (\autoref{sec:gespeicherte-ergebnisse}).
  \item[Befehl]\index{Befehl} Ein Zeichen- oder Modellierwerkzeug (\cmd{Line},
    \cmd{Box}, \cmd{Move} \dots), das aus der Werkzeugspalte, per Tastenkürzel oder durch
    Eintippen in die Befehlszeile gestartet wird (\autoref{kap:modellieren},
    \autoref{kap:befehle}).
  \item[Ansicht (Viewport)]\index{Ansicht} Eines der vier Zeichenfenster in der Mitte des
    Hauptfensters (Front, Left, Top, Perspective). Die \emph{aktive} Ansicht ist die zuletzt
    angeklickte (\autoref{sec:ansichten}).
\end{description}

\section{Aufbau des Handbuchs}\label{sec:aufbau}

Das Handbuch hat drei Teile.

\paragraph{Teil I -- Grundlagen.} Er führt durch die Bedienung des CAD-Programms ohne
Plugins:
\begin{itemize}
  \item \autoref{kap:installation} -- Installation auf allen drei Systemen, Einrichten der
    Python-Umgebung, erster Start, Updates, Absturzberichte, Einstellungen.
  \item \autoref{kap:hauptfenster} -- alle Bereiche des Hauptfensters und alle Menüs.
  \item \autoref{kap:navigation} -- Drehen, Verschieben, Zoomen, Auswählen, Ein- und
    Ausblenden.
  \item \autoref{kap:modellieren} -- Zeichnen und Modellieren mit Befehlen, Fang und
    Koordinateneingabe.
  \item \autoref{kap:organisieren} -- Ebenen, Gruppen, Modellbaum, Eigenschaften, Attribute,
    Rückgängig machen.
  \item \autoref{kap:dateien} -- Dokumente, STEP und STL, Beispielmodelle.
\end{itemize}

\paragraph{Teil II -- Berechnen mit Plugins.} \autoref{kap:plugins} erklärt, was alle
Plugins gemeinsam haben; \autoref{kap:workflow} den Knoten-Workflow. Danach folgt je ein
Kapitel für den Raytracer (\autoref{kap:raytracer}), den Path Tracer
(\autoref{kap:pathtracer}), FourierGrating (\autoref{kap:fouriergrating}) und Automotive
Lighting (\autoref{kap:automotive}) -- jeweils mit Beispielen zum Nachvollziehen.

\paragraph{Teil III -- Anhang.} Nachschlageteil: alle Tastenkürzel
(\autoref{kap:tastenkuerzel}), alle Befehle der Befehlszeile (\autoref{kap:befehle}), die
Programme und Pakete (\autoref{kap:programme}), alle Parameter und Aktionen jedes Plugins
(\autoref{kap:plugin-raytracer} bis \autoref{kap:plugin-automotive}), alle Blöcke des
Workflows (\autoref{kap:workflow-bloecke}), die Verwaltung von Plugins
(\autoref{kap:plugin-verwaltung}), das Dateiformat (\autoref{kap:dateiformat}), die
Fehlersuche (\autoref{kap:fehlersuche}) und das Glossar (\autoref{kap:glossar}).

Wer das Programm zum ersten Mal benutzt, liest am besten \autoref{kap:installation} bis
\autoref{kap:navigation} der Reihe nach und probiert die Schritte dabei aus. Die übrigen
Kapitel lassen sich unabhängig voneinander lesen.

\section{Schreibweisen}\label{sec:schreibweisen}

Die Oberfläche des Programms ist englisch. Menüs, Knöpfe und Felder stehen im Handbuch
deshalb so, wie sie auf dem Bildschirm erscheinen; erklärt wird auf Deutsch.

\begin{referenz}{>{\raggedright\arraybackslash}p{4.6cm}X}
\toprule
\textbf{Schreibweise} & \textbf{Bedeutung} \\
\midrule
\endhead
\menu{File,Save As...} & Ein Menübefehl: Menü \gui{File}, darin der Eintrag \gui{Save
  As...}. Untermenüs werden ebenso aneinandergereiht. \\
\gui{Check environment} & Beschriftung eines Bedienelements: Knopf, Feld, Kontrollkästchen,
  Dock, Fenster oder Spaltenkopf. \\
\keys{Strg+S} & Tastenkombination: Die Tasten werden gleichzeitig gedrückt, hier
  \keys{Strg} und \keys{S}. \\
\keys{Z}, dann \keys{E} & Tastenfolge: erst \keys{Z} drücken und loslassen, dann \keys{E}. \\
\maus{rechts ziehen} & Mausgeste: hier mit gedrückter rechter Maustaste ziehen. \\
\cmd{Line} & Ein Befehl, wie er in die Befehlszeile eingetippt werden kann. \\
\param{rays} & Name eines Plugin-Parameters oder eines Schlüssels in einer Datei. \\
\datei{examples/demo-parts.step} & Datei- oder Ordnername. \\
\code{x,y,z} & Eine Eingabe, wie sie getippt wird. \\
\plugin{Raytracer} & Name eines Plugins. \\
\block{Camera} & Block des Knoten-Workflows. \\
\bottomrule
\end{referenz}

\subsection*{Tasten auf dem Mac}
\index{Mac-Tasten}\index{Cmd-Taste}

Das Handbuch nennt Tastenkürzel in der Schreibweise von Windows und Linux. Auf dem Mac gilt:

\begin{itemize}
  \item \keys{Strg} entspricht der Befehlstaste \keys{Cmd} (\enquote{\textit{Command}},
    Kleeblatt-Symbol). \keys{Strg+S} heißt auf dem Mac also \keys{Cmd+S}, \keys{Strg+Shift+Z}
    heißt \keys{Cmd+Shift+Z}. Auch beim Auswählen mit der Maus (\keys{Strg}+Klick) ist auf dem
    Mac \keys{Cmd} gemeint.
  \item \keys{Alt} ist auf dem Mac die Wahltaste (\enquote{\textit{Option}}).
  \item \keys{Shift} ist die Umschalttaste, \keys{Enter} die Eingabetaste (Return),
    \keys{Entf} die Löschtaste nach rechts (\enquote{Delete}); \keys{Rück} ist die
    Rückschritttaste (auf dem Mac mit \enquote{delete} beschriftet). Beide löschen die
    Auswahl.
\end{itemize}

In den Menüs zeigt das Programm die Tastenkürzel in der Schreibweise des jeweiligen Systems
an, auf dem Mac also mit den Symbolen für Cmd, Shift und Option.

\begin{tipp}
Die Tastenkürzel ohne Zusatztaste (\keys{H}, \keys{Leertaste}, \keys{Z}, dann \keys{E} \dots)
wirken, solange eine Ansicht oder ein Dock den Fokus hat, nicht aber, während Sie in ein
Eingabefeld oder die Befehlszeile schreiben. Alle anderen Buchstaben, die Sie in einer Ansicht
tippen, landen in der Befehlszeile (\autoref{sec:befehlszeile}) -- auch \keys{H}, solange nichts
ausgewählt ist.
\end{tipp}

\subsection*{Kästen}

Im Text stehen drei Arten von Kästen:

\begin{hinweis}
Ergänzende Informationen, die man kennen sollte, um ein Verhalten zu verstehen.
\end{hinweis}

\begin{tipp}
Abkürzungen und Arbeitsweisen, die Zeit sparen.
\end{tipp}

\begin{warnung}
Stellen, an denen Daten verloren gehen können oder ein Ergebnis falsch werden kann.
\end{warnung}

\subsection*{Maßeinheiten}

Alle Längen im CAD-Modell sind Millimeter; das Dokumentformat speichert die Einheit
\code{mm}. Plugins können für ihre eigenen Parameter andere Einheiten verwenden (etwa
Mikrometer für Wellenlängen); sie stehen dann bei dem jeweiligen Parameter.
